\par\medskip
\noindent\underline{Corolario 3.5.} \underline{Sea $P$ un esquema en
grupos liso de presentación finita sobre $S$, con fibras conexas, y
sea}
\[
 0\longrightarrow Q'\longrightarrow Q\longrightarrow Q''
 \longrightarrow0
\]
\underline{una sucesión exacta de grupos conmutativos sobre $S$, con
$Q'$ un toro. Entonces el funtor de imagen inversa sobre las
biextensiones}
\begin{equation}\tag{3.5.1}
 \operatorname{BIEXT}(P,Q'';\underline G_m)
 \xrightarrow{\ \sim\ }
 \operatorname{BIEXT}(P,Q;\underline G_m)
\end{equation}
\underline{es una equivalencia de categorías. En particular, los
homomorfismos naturales siguientes son isomorfismos:}
\begin{equation}\tag{3.5.2}
 \Biext^i(P,Q'';\underline G_m)
 \xrightarrow{\ \sim\ }
 \Biext^i(P,Q;\underline G_m),
 \qquad i=0,1.
\end{equation}

En efecto, salvo equivalencia,
$\operatorname{BIEXT}(P,Q'';\underline G_m)$ se identifica con la
categoría de las biextensiones $E$ de $(P,Q)$ por $\underline G_m$
\underline{provistas} de una trivialización de la biextensión inducida
$E'$ de $(P,Q')$ por $\underline G_m$ (VII 3.7.6). Por 3.4 existe
exactamente una trivialización de $E'$ para cada biextensión $E$ dada.

%%%%%%%%%%  scan idx 274  |  folio 263  %%%%%%%%%%

Naturalmente, también puede demostrarse directamente que (3.5.2) es
un isomorfismo para $i=0,1$ —lo que equivale a la conclusión de 3.5—
utilizando la sucesión exacta de los $\Biext^i$ de VII 3.7.5 asociada
a
\[
 0\longrightarrow Q'\longrightarrow Q\longrightarrow Q''
 \longrightarrow0.
\]

\medskip
\noindent\underline{Observación 3.6.} El enunciado simétrico que se
obtiene de 3.5 intercambiando izquierda y derecha es también válido;
se deduce formalmente de 3.5 pasando a las biextensiones simétricas
(VII 2.7). Supongamos, además de los datos e hipótesis de 3.5, que se
tienen los datos e hipótesis simétricos: a saber, que $Q$ tiene fibras
conexas (equivalentemente, que $Q''$ tiene fibras conexas, pues $Q'$
las tiene), y que existe una sucesión exacta
\[
 0\longrightarrow P'\longrightarrow P\longrightarrow P''
 \longrightarrow0
\]
con $P'$ un toro. Aplicando sucesivamente 3.5 y su forma simétrica a
esta última sucesión y a $Q''$, se ve que el funtor
\begin{equation}\tag{3.6.1}
 \operatorname{BIEXT}(P'',Q'';\underline G_m)
 \xrightarrow{\ \sim\ }
 \operatorname{BIEXT}(P,Q;\underline G_m)
\end{equation}
es una equivalencia de categorías. En particular, el homomorfismo
siguiente es un isomorfismo:
\begin{equation}\tag{3.6.2}
 \Biext^i(P'',Q'';\underline G_m)
 \xrightarrow{\ \sim\ }
 \Biext^{\sisi{1}{i}}(P,Q;\underline G_m),
 \qquad i=0,1.
\end{equation}

\medskip
\noindent\underline{Proposición 3.7.} \underline{Sea $A$ un esquema
abeliano sobre $S$, sea $T$ un esquema en grupos de tipo multiplicativo
(SGA 3 IX) y de tipo finito sobre $S$, y sea}
\[
 \underline M=\uHom(T,\underline G_m)
\]
\underline{el dual de Cartier de $T$ (un grupo constante retorcido
finitamente generado sobre $S$, por 3.3.3). Entonces las categorías
$\operatorname{EXT}(A,T)$ y
$\operatorname{BIEXT}(A,\underline M;\underline G_m)$, formadas
respectivamente por las extensiones conmutativas de $A$ por $T$ y por
las biextensiones de $(A,\underline M)$ por $\underline G_m$, son
rígidas y canónicamente equivalentes. Sus grupos de clases de
isomorfismo son canónicamente isomorfos a $\Hom(\underline M,A^*)$,
donde}

%%%%%%%%%%  scan idx 275  |  folio 264  %%%%%%%%%%

\[
 A^*=\uExt^1(A,\underline G_m)
\]
\underline{es el esquema abeliano dual de $A$ (3.2). En particular,
se tiene un isomorfismo canónico}
\begin{equation}\tag{3.7.1}
 \Ext^1(A,T)\xrightarrow{\ \sim\ }
 \Hom(\underline M,A^*).
\end{equation}

Este es el caso particular de 1.6 que se obtiene tomando $P=A$ y
$Q=\underline M$. Por el teorema de bidualidad (SGA 3 X 5.7), se
tiene en efecto
\[
 T\simeq\uHom(\underline M,\underline G_m),
\]
y solo queda verificar las condiciones (1.6.1), que aquí se escriben
\[
 \uHom(A,\underline G_m)=0,
 \qquad
 \Hom\bigl(A,\uExt^1(\underline M,\underline G_m)\bigr)=0.
\]
La primera igualdad se recordó en (3.2.1); la segunda puede reforzarse
como
\[
 \uExt^1(\underline M,\underline G_m)=0.
\]
Esta afirmación es local sobre $S$, por lo que podemos suponer que
$\underline M$ es constante, $\underline M=M_S$. Descomponiendo $M$
como suma directa de $\mathbb Z$-módulos cíclicos, quedamos reducidos a
$M=\mathbb Z$, que es trivial, y a
$\sisi{\mathbb Z}{M}=\mathbb Z/n\mathbb Z$, que se deduce de la
$n$-divisibilidad de $\underline G_m$ (o, más eruditamente, de 3.3.1).

Cuando $T=\underline G_m$, y por tanto
$\underline M=\mathbb Z_S$, (3.7.1) se reduce al isomorfismo identidad
(3.2.3). Naturalmente, una demostración directa de (3.7.1) sería
esencialmente trivial, sin introducir el grupo
\[
 \Ext^1\bigl(A\mathbin{\mathop{\otimes}^{\mathbb L}}\underline M,
              \underline G_m\bigr)
 =\Biext^1(A,\underline M;\underline G_m).
\]
